% ============================================================
% Chapter 4: Proposed Method: Specification and Design
%   Section: Concurrency Design
% Included from chapters/04-proposed-method/proposed-method.tex
% Source: docs/decisions/concurrency.md; src/services/concurrency.py.
% ============================================================

\section{Concurrency Design}
\label{sec:design-concurrency}

Checking one article used to be sequential end to end: per claim, the
searches one after another, the page fetches one after another, some
twenty embedding calls one at a time and one model call, then the next
claim. A four-claim article spent minutes waiting, and none of that
waiting was on the processor. Three changes, in the order they were
worth making, removed most of it: the embedding calls were batched; a
claim's queries and page fetches, and the web search and the corpus
search, were run in parallel; and claims were run in parallel with one
another.

The design rests on four rules.

\begin{enumerate}
    \item \textbf{A claim's own stages stay sequential.} Retrieval feeds
    ranking, ranking feeds the model, and the model's answer is what
    recalibration works on. The concurrency is between claims, never
    inside one.
    \item \textbf{Limits belong to the resource.} Fan-out multiplies:
    three articles times four claims times eight page fetches is 96
    simultaneous requests from one machine, enough for the search
    engines behind SearXNG to start refusing it. The real ceilings are
    therefore one semaphore per external service --- two concurrent
    searches, one inference call, two model calls by default (one on a
    CPU), eight page fetches, two browser pages --- applied inside the
    client that talks to each service. They are process settings, not
    per-run thresholds: a per-request override would let one caller
    decide how hard the process leans on a service every other caller
    shares.
    \item \textbf{A permit is held only around a single call.} Holding
    one while waiting on work that needs the same resource is the
    bounded-pool deadlock, in which every worker holds what every other
    is waiting for.
    \item \textbf{Results come back in input order.} The evidence list's
    positions are what the model cites by number and what the response,
    the cache and the lake record. Results in completion order would
    silently point every citation at a different source --- a wrong
    answer, not a crash. A test pins it.
\end{enumerate}

Threads were chosen over asynchronous code because everything the
pipeline calls --- the HTTP clients, trafilatura, the embedded vector
store, the model client --- is synchronous, and the workload is waiting
on input and output, for which threads are adequate. Progress events
from several claim threads are serialised in one place, so the code that
receives them keeps a single-threaded contract, and each event carries
the index of its claim, so a client can follow claims whose events
interleave.

The speed-up was measured against a simulated network at the declared
ceilings: $3.4\times$ on a four-claim article with a fast model, and
$1.16\times$ with the production CPU model at one call at a time, where
the model takes 180 of the 184 seconds and nothing else is left to
overlap. The simulation's 184~s matches the production stack's real run
of the same article.
